Another thing we can do to improve our search performance comes from the \textbf{heavy tail behaviour}. Assuming a Constraint Satisfaction Problem that 
requires some inputs (\textit{e.g. in the n-queens problem the number of queens is an input parameter}), its complexity depends from the size of 
the input parameters, even though the model still always the same. \vspace{7pt}

We often observe some exceptionally hard instances that take exceptionally longer time to solve. Of course, solving an instance of the n-queens problem 
within $1000$ of queens takes much more time than solve it with $4$ queens. \vspace{7pt}

The thing is: \textbf{why this is happening}? \\
The main observation is that: this occurs not because any instance of the problem is itself difficult to solve by nature, but because we are playing
with different parameters of the solver. For some combinations between the instances of the problem and the solver parameters, we get trapped into an 
exponential subtree, making the problem difficult to solve. \vspace{7pt}

At the beginning of the search process all the heuristics used are unreliable: if we make a mistake early during search, we get stuck in a subtree.
Therefore, the conclusion is that we don't really know how heuristics work, such mistakes or wrong decisions seem to be totally \textbf{random}. \vspace{7pt}

In order to avoid these random mistakes we can do:
\begin{itemize}
    \renewcommand{\labelitemi}{-}
    \item \textbf{Randomization}. \\ 
    Since these mistakes are random, why don't we introduce \textbf{randomization} in our solvers also? By randomization we can kill this effect associated to 
    randomness. For instance we could: 
    \begin{enumerate}
        \item Pick some variables and/or values randomly.
        \item Break ties randomly.
    \end{enumerate}
    \item \textbf{Restarting}. \\
    \textbf{Restart} the search process can be very useful when the algorithm keeps backtracking without returning us anything. For sure, to be effective 
    we have to search in a different way anytime we restart the process. Searching anytime differently we also introduce:
    \begin{enumerate}
        \item Randomization.
        \item Learning from the accumulated experiences of previous runs, or in other words previous runs can be effective for the following ones.
    \end{enumerate}
    \item \textbf{Randomization + Restarting}. \\
    We could also combine together the previous strategies, deleting the huge variance in solver performance. 
\end{itemize}